\documentclass[11pt,a4paper]{article}
\usepackage[margin=27mm]{geometry}
\usepackage{amsmath,amssymb,amsthm,booktabs,longtable,array}
\usepackage{fontspec,unicode-math}
\setmainfont{STIXGeneral}
\setmathfont{STIX Two Math}
\usepackage{hyperref}
\setlength{\emergencystretch}{2em}
\newtheorem{theorem}{Theorem}
\newtheorem{lemma}[theorem]{Lemma}
\newtheorem{corollary}[theorem]{Corollary}
\newtheorem{proposition}[theorem]{Proposition}
\newcommand{\PSL}{\operatorname{PSL}}
\newcommand{\PGL}{\operatorname{PGL}}
\newcommand{\tr}{\operatorname{tr}}
\newcommand{\Aut}{\operatorname{Aut}}
\newcommand{\ip}{\mathrm{IP}}
\title{Common theory for projective markings and rotations}
\author{The human\_compression project}\date{11 October 2026}
\begin{document}\maketitle
\begin{abstract}
Projective centralizers and matrix intertwiners control both involutory
string markings and rotational markings. We retain the complete common
proof for the three existing polytope and maniplex applications, then
prove simultaneous reversal, a one-parameter classification of two affine
ends, and one trace quadratic for the three exceptional rank-four families.
An anchored-involution interface supplies the chiral rank bound; the
previously accepted norm-polynomial and translated-subfield arguments
supply the odd-composite-degree construction. Finite presentations and
small-field exceptions are paid by complete certificates. Classical
subgroup, automorphism, character-sum and geometric correspondence inputs
remain explicit. Five source applications use this whole common paper.
\end{abstract}
\section{Markings and retained inputs}
A \emph{string marking} of a group $G$ is a generating ordered tuple
$(r_0,\ldots,r_{n-1})$ of distinct nonidentity involutions with
$r_ir_j=r_jr_i$ whenever $|i-j|>1$. Put
$G_I=\langle r_i:i\in I\rangle$, with $G_\varnothing=1$.
Its labels are $|r_0r_1|,\ldots,|r_{n-2}r_{n-1}|$.
The marking is a string C-marking if, in addition,
\begin{equation}\label{eq:ip}
 G_I\cap G_J=G_{I\cap J}\qquad(I,J\subseteq\{0,\ldots,n-1\}).
\end{equation}
For rank four write $H=\langle r_0,r_1,r_2\rangle$ and
$J=\langle r_1,r_2,r_3\rangle$ for its two ends. Throughout,
$L(q)=\PSL_2(q)$ and $P(q)=\PGL_2(q)$, for a prime power $q$.

We retain the following classical inputs, in exactly the forms used below.
Dickson's classification says that a proper subgroup of $L(q)$ or $P(q)$,
for odd $q$, is cyclic, dihedral, elementary abelian of characteristic order,
a subgroup of an affine group $V\rtimes k^\times$, one of $A_4,S_4,A_5$,
or a conjugate of a standard subfield group $L(k)$ or $P(k)$.
Only groups which actually embed are intended. In the subfield alternatives
$k\subseteq\mathbb F_q$ and the embedding is the natural matrix embedding,
not an arbitrary abstract group isomorphism. The possible same-field proper
subgroup $L(q)<P(q)$ must be retained. We use
\[
 |L(q)|=\frac{q(q^2-1)}{(2,q-1)},\qquad |P(q)|=q(q^2-1),
\]
the simplicity of $L(q)$ for $q\ge4$, and $C_{P(q)}(L(q))=1$.
The small identifications are $L(2)=P(2)\cong S_3$, $L(3)\cong A_4$,
$P(3)\cong S_4$, $L(5)\cong A_5$ and $P(5)\cong S_5$.
These subgroup and structural facts are classical; subgroup multiplicity
formulas from the sources are not inputs.

We also retain the correspondence between ordered string C-markings and
regular abstract polytopes, including the following clauses. Rank-$i$ faces
are cosets of $G_{\{0,\ldots,n-1\}\setminus\{i\}}$; reversal gives duality;
a marking-preserving epimorphism of string C-groups induces a covering.
The local markings $[3,5]_5$ and $[5,3]_5$ below are the usual antipodal
icosahedral and dodecahedral quotients. These standard local-map names and
the simplex interpretation are retained geometric identifications.
They do not replace the required verification of the actual markings.

All groups $L(q),P(q)$ are centre-free and directly indecomposable. For the
simple groups this is immediate. To see the assertion for $P(q)$ when
$q\ge4$, let commuting subgroups $A,B$ generate $P(q)$; both are normal.
The subgroup $[L(q),A]$ lies in $L(q)\cap A$ and is normal in the simple
group $L(q)$. If it is nontrivial, $L(q)\le A$, so $B$ centralizes $L(q)$
and is trivial. Otherwise $A$ centralizes $L(q)$ and is trivial. The small
cases follow from the normal-subgroup lists
$1,A_3,S_3$; $1,V_4,A_4$; and $1,V_4,A_4,S_4$.
In any string marking of these groups every adjacent product has order
at least three. Indeed, a commuting adjacent pair splits the generators
into two nontrivial commuting generated subgroups. Their intersection is
central, hence trivial, and the resulting direct product contradicts
indecomposability. This argument does not assume \eqref{eq:ip}.

\section{Common centralizers}
\begin{lemma}\label{lemma:centralizer}
In characteristic two, the centralizer of a nonidentity involution in
$P(q)$ is abelian. Let $k\subseteq K$ be finite fields of odd characteristic,
and let $a,b\in P(k)$ be involutions with $|ab|>2$. Their common centralizer
in $P(K)$ lies in $P(k)$. It is trivial when $ab$ is unipotent and has
order two otherwise. In the latter case its nonidentity element is represented
by the invertible matrix
\begin{equation}\label{eq:trace}
 Z=A-\tfrac12\tr(A)I,
\end{equation}
where $A\in\operatorname{GL}_2(k)$ represents $ab$.
\end{lemma}
\begin{proof}
In characteristic two an involution is unipotent. Over an algebraic closure
conjugate it to $x\mapsto x+1$. Its unique fixed point is infinity, so a
commuting projectivity is $x\mapsto cx+d$. Commutation gives $c=1$;
these translations form an abelian group. The centralizer over the original
field is a subgroup of this group.

In odd characteristic first suppose $u=ab$ is semisimple. Over the algebraic
closure write $u:x\mapsto\lambda x$, with $\lambda\ne\pm1$.
A commuting projectivity preserves each of its two fixed points, since
interchanging them would send $u$ to $u^{-1}\ne u$. It is therefore
$x\mapsto cx$. As $aua^{-1}=u^{-1}$, the involution $a$ has the form
$x\mapsto d/x$. Commutation with $a$ gives $c^2=1$, so the common
centralizer is $\{1,x\mapsto-x\}$. If $A$ has eigenvalues
$\alpha,\beta$, the matrix in \eqref{eq:trace} has the nonzero eigenvalues
$(\alpha-\beta)/2$ and $(\beta-\alpha)/2$. It represents $x\mapsto-x$
and has entries in $k$. Rescaling $A$ only rescales $Z$, proving the
claimed descent, including the nonsplit case.

If $u$ is unipotent, conjugate it to $x\mapsto x+1$. Its centralizer again
consists of translations $x\mapsto x+c$. An involution conjugating it to
its inverse has the form $x\mapsto-x+d$. A commuting translation satisfies
$2c=0$, and hence $c=0$. These cases exhaust $u$.
\end{proof}

\begin{lemma}\label{lemma:rank}
Let $G$ be centre-free with a string marking. Its first two generators do
not commute. If every noncommuting pair of involutions has at most one
nonidentity common centralizer element, the rank is at most four. If every
involution centralizer is abelian, the rank is at most three.
\end{lemma}
\begin{proof}
If $r_0$ commuted with $r_1$, it would commute with every generator by the
string relations, so would be a nonidentity central element. At rank at
least five, both $r_3$ and $r_4$ centralize $r_0,r_1$, contradicting the
first centralizer hypothesis and their distinctness. Under the second
hypothesis a rank-four marking puts $r_0,r_1$ in the abelian centralizer
of $r_3$, again a contradiction.
\end{proof}

\begin{lemma}\label{lemma:odd}
Let a rank-four string marking lie in $P(q)$, with $q$ odd and adjacent
labels greater than two. If $H$ is proper, then $|r_0r_1|$ is odd; if
$J$ is proper, then $|r_2r_3|$ is odd.
\end{lemma}
\begin{proof}
The nonidentity element $r_3$ centralizes $D=\langle r_0,r_1\rangle$.
Lemma~\ref{lemma:centralizer} excludes the unipotent case and gives a
unique nonidentity common centralizer element. If $m=|r_0r_1|$ were even,
$(r_0r_1)^{m/2}$ would be that element. Then $r_3\in H$, contrary to
$G=\langle H,r_3\rangle$ and properness. Reverse the indices for $J$.
\end{proof}

\begin{theorem}[Projective reduction]\label{thm:projective}
A string marking of $L(q)$ or $P(q)$ has rank at most four; rank four
requires odd $q$. In rank four every proper end of $L(q)$ is $S_4$ or
$A_5$. Every proper end of $P(q)$ is $S_4,A_5$, or its natural subgroup
$L(q)$. If both ends are proper, both are $S_4$ or $A_5$ and both outer
labels are odd. These statements require no intersection property.
\end{theorem}
\begin{proof}
The rank claims follow from Lemmas~\ref{lemma:centralizer} and
\ref{lemma:rank}, since the projective groups are centre-free. This also
handles $q=2$; $L(3)$ has no generating tuple of involutions, since its
involutions generate only $V_4$. All adjacent labels in the remaining
cases exceed two by the first section.

Let $H=\langle r_0,r_1,r_2\rangle$ be proper and put
$D=\langle r_0,r_1\rangle$. A cyclic or elementary abelian group cannot
contain the nonabelian group $D$. Involutions in $A_4$ generate only
$V_4$. In an affine group in odd characteristic every involution is
$x\mapsto-x+v$; two such involutions commute only when they coincide.
Thus the distinct commuting endpoints $r_0,r_2$ exclude this alternative,
including affine dihedral groups of order $2p$.

If $H$ is dihedral, its distinct commuting involutions $r_0,r_2$ force
its rotation subgroup to have even order. (In an odd dihedral group two
reflections commute only when equal.) Its central involution $z$ and
$r_3$ both centralize $D$. Lemma~\ref{lemma:centralizer} gives $r_3=z$,
so $r_3\in H$, a contradiction. This exclusion uses properness, not an
intersection identity inside $H$.

If $H$ is a natural subfield group, conjugate it into $P(k)$.
Lemma~\ref{lemma:centralizer} puts $r_3$ in the same $P(k)$, hence puts
all of $G$ there. For a proper subfield of size $s$ we have $q\ge s^2$,
and
\[
 |G|\ge q(q^2-1)/2>s(s^2-1)=|P(k)|,
\]
a contradiction. The only remaining same-field proper alternative is
$H=L(q)<P(q)$. This proves the asserted lists of proper ends. The
qualification ``natural'' is essential: exceptional $A_5\cong L(5)$
and $S_4\cong P(3)$ in other characteristics are not subfield exclusions.

Suppose now that both ends are proper. Lemma~\ref{lemma:odd} makes both
outer labels odd. If $H=L(q)<P(q)$, then $H$ is normal of index two;
$r_2\in H$ and $r_3\notin H$. The image of $r_2r_3$ in $P(q)/H$
therefore has order two, impossible for its odd order. Reversal excludes
the same alternative for $J$. The remaining ends are $S_4,A_5$.
\end{proof}

\section{One finite theorem for exceptional ends}
For a triple $f=(m,n,k)$ write
\begin{equation}\label{eq:local}
 [m,n]_k=\langle x,y,z\mid x^2=y^2=z^2=(xz)^2=
 (xy)^m=(yz)^n=(xyz)^k=1\rangle.
\end{equation}
The third parameter is the Petrie order, not an adjacent label.
The six ordered local classes are
\begin{equation}\label{eq:localtypes}
\begin{array}{c|ccc}
 S_4&(3,3,4)&(3,4,3)&(4,3,3)\\
 A_5&(3,5,5)&(5,3,5)&(5,5,3).
\end{array}
\end{equation}
For $f=(m,n,k)$ and $v=(n,s,l)$ in this list define
\begin{equation}\label{eq:universal}
\begin{split}
 U(f,v)=\langle r_0,r_1,r_2,r_3\mid{}
 &r_i^2=1,\ (r_ir_j)^2=1\ (j-i>1),\\
 &(r_0r_1)^m=(r_1r_2)^n=(r_2r_3)^s=1,\\
 &(r_0r_1r_2)^k=(r_1r_2r_3)^l=1\rangle.
\end{split}
\end{equation}
Matching the middle label gives fourteen ordered pairs. Reversal sends
$(f,v)$ to $(v^*,f^*)$, where $(a,b,c)^*=(b,a,c)$.
There are ten reversal orbits, six fixed. Their complete table is:
\begin{center}\small
\begin{tabular}{rcccccl}
\toprule
Row&$f$&$v$&Order&Group&C-marking&Outer labels\\
\midrule
1&$(5,5,3)$&$(5,5,3)$&1&trivial&no&$5,5$\\
2&$(5,5,3)$&$(5,3,5)$&1&trivial&no&$5,3$\\
3&$(5,3,5)$&$(3,5,5)$&3420&$L(19)$&yes&$5,5$\\
4&$(5,3,5)$&$(3,4,3)$&60&$A_5$&no&$5,4$\\
5&$(5,3,5)$&$(3,3,4)$&1&trivial&no&$5,3$\\
6&$(4,3,3)$&$(3,4,3)$&96&$2^4:S_3$&yes&$4,4$\\
7&$(4,3,3)$&$(3,3,4)$&24&$S_4$&no&$4,3$\\
8&$(3,5,5)$&$(5,3,5)$&660&$L(11)$&yes&$3,3$\\
9&$(3,4,3)$&$(4,3,3)$&1&trivial&no&$3,3$\\
10&$(3,3,4)$&$(3,3,4)$&120&$S_5$&yes&$3,3$\\
\bottomrule
\end{tabular}
\end{center}
Here an entry ``yes'' means that the specified generators, not merely some
marking of the abstract group, form a string C-group. Rows which collapse
do not have the indicated local groups; the local triples label their
presentation input. The order-$96$ row has adjacent outer labels $4,4$,
not $3,3$.

\begin{lemma}[Exact finite certificate]\label{lemma:finite}
Every rank-three string marking of $S_4$ or $A_5$ has the intersection
property. Up to the full automorphism group its ordered type and Petrie
order are exactly \eqref{eq:localtypes}; \eqref{eq:local} is an exact
presentation in each case. The ten-row table above is correct, including
its groups, specified markings, collapse assertions and reversal coverage.
In rows $3,8,10$, reversal extends to an automorphism of the marked group.
\end{lemma}
\begin{proof}
The complete program \texttt{required\_finite.g} enumerates all involutions
in $S_4,A_5$, all ordered pairs of distinct commuting endpoints $x,z$,
and every middle involution $y$. Its admission tests are only distinctness
and generation. It admits $72$ and $360$ triples respectively. For every
admitted triple it then checks both adjacent orders and all $64$
subset-pair identities in \eqref{eq:ip}. Thus IP is an output, never a
filter whose use would assume the desired category change. The full
automorphism orbits give the three classes in each row of
\eqref{eq:localtypes}. The code constructs each literal presentation
\eqref{eq:local} and computes its order $24$ or $60$. The marking
provides a surjection onto the corresponding finite group; equality of
orders makes it an isomorphism.

The code independently generates all fourteen compatible pairs, partitions
by reversal and checks the ten supplied representatives and six fixed
pairs. For each literal presentation \eqref{eq:universal}, it computes
the finite presentation order, constructs an isomorphic permutation group
and transports the four original generators. It checks generator orders and distinctness, then tests IP over all sixteen
subsets. Every accepted marking satisfies all $256$ subset-pair identities;
a rejected row may fail a generator condition or terminate at its first
failed intersection. Thus the negative entries certify failure of the
complete string C-marking conditions, not necessarily a separate IP test.
Rows $3,8,10$ are identified by actual group isomorphisms with
$L(19),L(11),S_5$; generator reversal is checked by constructing the
homomorphism and testing bijectivity. In rows $4,7$ it verifies the indicated group and that
one end is the whole group. In row $6$ it verifies an elementary abelian
normal subgroup of order $16$, quotient $S_3$ and a complement. The four
trivial rows are checked directly. These exact computations prove the
finite assertions relative to GAP's retained finite-group and finite-
presentation algorithms; they are not claimed as Lean kernel computations.
\end{proof}

\begin{theorem}[Exceptional ends and odd outer labels]\label{thm:finite}
Let $(G;r_0,r_1,r_2,r_3)$ be a string marking whose two ends are proper
and isomorphic to $S_4$ or $A_5$. If both outer labels are odd, then,
up to isomorphism preserving the ordered marking, it is exactly one of
rows $3,8,10$. Each row supplies one class, satisfies IP, and has type
$\{5,3,5\}$, $\{3,5,3\}$ or $\{3,3,3\}$ respectively.
In particular the statement does not require $G$ to be a projective group.
\end{theorem}
\begin{proof}
The local census and its exact presentations give a marking-preserving
surjection from one of the ten groups $U(f,v)$ to $G$, after possibly
reversing the marking. It is necessary to examine quotients, not just
the universal groups themselves. The trivial rows cannot give nonidentity
markers. In rows $4,7$, an end already equals the full universal group;
this equality survives every quotient, contradicting properness. Row $6$
has first and last local product orders four, whereas the given ends
have their exact local markings and hence these exact orders. It is
excluded by oddness.

The remaining universals are the simple groups $L(19),L(11)$ and $S_5$.
The simple groups have no nontrivial proper quotient. The normal subgroups
of $S_5$ are $1,A_5,S_5$, so its only nontrivial proper quotient is $C_2$,
which cannot have an end $S_4$ or $A_5$. Thus every surviving map is an
isomorphism on all four distinguished generators. Lemma~\ref{lemma:finite}
gives IP and the asserted types. The local classes were taken under full
automorphisms, so no marking was lost through a conjugacy restriction.
The reversal automorphisms verified below identify the two representatives
of each reversal orbit. There is consequently one ordered-marking class
per surviving row.
\end{proof}

\noindent\emph{Sharpness of the parity hypothesis.} The verified order-$96$
row has proper $S_4$ ends, type $\{4,3,4\}$ and IP. It therefore
contradicts the three-case conclusion if the odd-outer-label hypothesis
is removed. We make no classification of all quotients without that
hypothesis.

\begin{corollary}\label{cor:projective}
Every rank-at-least-four string marking of $L(q)$ or $P(q)$ with both
ends proper is one of these three markings. The first two occur in
$L(19),L(11)$ respectively; the third occurs in $P(5)$.
A string C-marking automatically has both ends proper, so the same
conclusion applies to all projective string C-markings of rank at least four.
\end{corollary}
\begin{proof}
Theorem~\ref{thm:projective} supplies rank four, the exceptional ends and
odd outer labels; apply Theorem~\ref{thm:finite}. The order formulas are
strictly increasing in odd $q$. The groups $L(11),L(19)$ are simple,
whereas odd-characteristic $P(q)$ has a proper subgroup of index two.
The group $S_5$ is not simple and has order $120=|P(5)|$; the small
$L(3)$ is not $S_5$. These observations give exactly the stated field
values and projective types. Finally, if an omitted generator belonged
to its complementary subgroup, \eqref{eq:ip} would intersect that subgroup
with its generator subgroup nontrivially. Hence C-markings have proper ends.
\end{proof}

\section{Geometry and universal coverings}
\begin{theorem}\label{thm:geometry}
The three markings in Theorem~\ref{thm:finite} give the 57-cell, the
11-cell and the 4-simplex. Their types, local maps and face vectors are
\[
\begin{array}{c|c|c|c}
\text{group}&\text{type}&\text{facet / vertex figure}&(f_0,f_1,f_2,f_3)\\
L(19)&\{5,3,5\}&\text{hemi-dodecahedron / hemi-icosahedron}&(57,171,171,57)\\
L(11)&\{3,5,3\}&\text{hemi-icosahedron / hemi-dodecahedron}&(11,55,55,11)\\
S_5&\{3,3,3\}&\text{tetrahedron / tetrahedron}&(5,10,10,5).
\end{array}
\]
All three are self-dual and universal among regular polytopes with their
respective facet and vertex-figure types. The 11-cell edge graph is
$K_{11}$. The even-word subgroup is the full group in the two simple
rows, and is $A_5$ in the simplex row.
\end{theorem}
\begin{proof}
The two simple rows have local presentations $[3,5]_5,[5,3]_5$, with the
retained antipodal map interpretations. Their exact global presentations
are the string Coxeter relations of types $\{3,5,3\}$ and $\{5,3,5\}$,
together with
\[
 (r_0r_1r_2)^5=(r_1r_2r_3)^5=1.
\]
For the third row send the four generators to $(12),(23),(34),(45)$.
These generate $S_5$ and satisfy every relator in
\eqref{eq:universal}, including the two fourth powers of consecutive
triple products. Both groups have order $120$, so this is an isomorphism
of marked groups with the standard simplex marking.

The program verifies that generator reversal extends to an automorphism
in each of the three rows; the retained correspondence gives self-duality.
It also computes each of the four complementary-subgroup indices, giving
the displayed face vectors. In the simple rows this has the following
direct interpretation. Both ends have order $60$. An edge stabilizer is
$\langle r_0\rangle\times\langle r_2,r_3\rangle$, since the factors
commute and IP makes their intersection trivial. Its order is $4m$ for
outer label $m=3$ or $5$; the two-face stabilizer has the same order.
Thus $660/60=11$, $660/12=55$, $3420/60=57$ and $3420/20=171$ give the
first two vectors. For the 11-cell the code constructs actual vertex and
edge cosets and checks their incidences: every edge has two endpoints and
the $55$ resulting unordered pairs are distinct. There are exactly $55$
pairs on eleven vertices, proving $K_{11}$ rather than inferring it merely
from face counts.

The even-word subgroup is normal of index at most two: in its quotient
all involutions have the same image. Simplicity makes it the full group
in the two simple rows. In the simplex marking parity is the ordinary
sign homomorphism, giving $A_5$.

Let $Q$ be any regular polytope with one of the three indicated pairs of
local maps. Its distinguished generators satisfy the corresponding
literal presentation \eqref{eq:universal}. There is therefore a marked
epimorphism from the already verified string C-group of that row to its
automorphism group. The retained covering theorem gives a covering of
$Q$ by the corresponding polytope above. No finiteness assumption on $Q$
is needed: the epimorphism itself forces its group to be finite. This is
universality of these exact markings, not an assumption that every
abstract amalgam has IP.
\end{proof}

\section{The common extension construction}
We use the standard torus structure in $L(q)$ for odd $q$: its split and
nonsplit cyclic tori have orders $(q-1)/2,(q+1)/2$, with dihedral
normalizers of twice these orders. The split torus consists of diagonal
projectivities of square eigenvalue ratio. The nonsplit torus comes from
multiplication on $\mathbb F_{q^2}$ modulo scalars, with the determinant-
square condition. Reflections in the normalizer interchange the two
eigenlines. These classical torus facts are retained; no maximality of
these normalizers is assumed.

\begin{theorem}[Half-turn extension]\label{thm:extension}
Let $q$ be odd, choose $\alpha\in\{1,-1\}$ with $q\equiv\alpha\pmod4$,
and put $N=(q-\alpha)/2$. Let $(a,b,c)$ be a rank-three string marking
of $H\le L(q)$, with $|bc|>2$ and $|ab|=m>1$ an odd divisor of $N$.
There is a unique nonidentity projectivity $z$ centralizing $a,b$;
it lies in $L(q)$ and $(a,b,c,z)$ is a string marking of $\langle H,z\rangle$.
If $H=L(q)$, this group remains $L(q)$. If $H$ is maximal proper in
$L(q)$ and $C_H(\langle a,b\rangle)=1$, it is again $L(q)$ and its
first end is exactly $H$.
\end{theorem}
\begin{proof}
The element $ab$ is semisimple, since $m$ divides $q-\alpha$ and hence
is prime to the characteristic. It has two eigenlines over at most a
quadratic extension. The two torus orders are coprime consecutive integers;
since $m>1$ divides $N$, its torus has order $N$. Both $a,b$ invert
$ab$, so preserve its unordered eigenline pair and normalize this torus.
The even cyclic torus has a unique involution $z$. It commutes with the
torus and its inverting normalizer coset, so centralizes $a,b$.
Lemma~\ref{lemma:centralizer} proves its asserted uniqueness even in
$P(q)$.

The dihedral group $\langle a,b\rangle$ has odd rotation order and
trivial centre. Thus $z$ does not belong to it and differs from $a,b$.
It also differs from $c$, because it commutes with $b$ whereas $|bc|>2$.
The four generators are distinct nonidentity involutions and have every
required distant commutation. They generate $\langle H,z\rangle$ by
construction. If $H=L(q)$ there is nothing further to prove. Under the
second hypothesis $z\notin H$, because its membership would give a
nonidentity element of the stated trivial internal centralizer.
Maximality then gives $\langle H,z\rangle=L(q)$.
\end{proof}
The extension theorem concerns string markings; it does not assert IP
for the extended object. The maniplex application proves its category
correspondence and hence the assertion that the first end describes the
original map. Its all-prime constructions and explicit source-basis
identifications are local proofs in that application, not finite-sample
conclusions.


\section{Rotations and their retained geometric interface}
For the new applications we retain the rotation-group correspondence of
Schulte--Weiss \cite{SW}. A rank-$n$ rotation marking is an ordered
generating tuple $(\sigma_1,\ldots,\sigma_{n-1})$ with every interval
product $\sigma_i\cdots\sigma_j$, $i<j$, a nonidentity involution.
For a set $I$ of vertex indices put
\[
 G_I^+=\langle\sigma_i\cdots\sigma_j:i\le j,\ i-1,j\in I\rangle.
\]
The full rotation intersection property is
$G_I^+\cap G_J^+=G_{I\cap J}^+$. In rank four it is equivalent to
the two adjacent cyclic intersections being trivial and the two ends
intersecting in the middle cyclic group. Such a marking gives a chiral
polytope exactly when it has no involutory automorphism sending
$\sigma_1$ to $\sigma_1^{-1}$, $\sigma_2$ to $\sigma_1^2\sigma_2$
and fixing the later rotations. Otherwise it gives a directly regular
polytope. This correspondence, including the equivalence of intersection
conditions, is retained rather than claimed as a new theorem.

We also retain the full semilinear automorphism theorem
$\operatorname{Aut}(L(q))=\operatorname{P\Gamma L}_2(q)$ for $q\ge4$
and its corresponding $P(q)$ form, with the standard small identifications.
Finite-field norm and cyclic-character facts and the exact multiplicative
character bound from \cite{CKR} are the only new analytic inputs.
The projective subgroup alternatives are the same classical inputs as
in the first section, used in their natural-subfield form.
\section{One simultaneous reversal for two projectivities}
\begin{theorem}\label{E.SimultaneousReversal}
Let $F$ be a field and let $A,B\in\mathrm{GL}_2(F)$ have no common
invariant line over an algebraic closure. Then $X=AB-BA$ is invertible,
$X^2$ is scalar, and conjugation by $[X]$ simultaneously inverts
$[A],[B]\in\mathrm{PGL}_2(F)$. In particular, every generating pair of
$\mathrm{PSL}_2(q)$ or $\mathrm{PGL}_2(q)$ admits an automorphism of
order dividing two which inverts both generators.
\end{theorem}
\begin{proof}
The Cayley--Hamilton identity gives
\[
 XA+AX=\operatorname{tr}(A)X,\qquad
 XB+BX=\operatorname{tr}(B)X.
\]
Indeed, the first left side is $A^2B-BA^2$, and the second is
$AB^2-B^2A$. Thus $XA=\operatorname{adj}(A)X$ and
$XB=\operatorname{adj}(B)X$. These identities hold in every
characteristic. Also $\operatorname{tr}(X)=0$, whence
$X^2=-\det(X)I$.

To see that $X$ is invertible, extend scalars to an algebraic closure.
Neither $A$ nor $B$ is scalar, since otherwise they share an eigenline.
If $A=\operatorname{diag}(a,d)$ with $a\ne d$ and
$B=\left(\begin{smallmatrix}u&v\\w&z\end{smallmatrix}\right)$,
then $\det(AB-BA)=(a-d)^2vw$. Its vanishing makes either coordinate
line common and invariant. Otherwise $A$ has one Jordan block and,
after rescaling its off-diagonal entry, the commutator is
$\left(\begin{smallmatrix}w&z-u\\0&-w\end{smallmatrix}\right)$.
Singularity again makes the unique eigenline of $A$ invariant under
$B$. Both possibilities contradict the hypothesis. Since
$\operatorname{adj}(A)=\det(A)A^{-1}$, the two matrix identities now
prove the projective assertion.

For the finite projective groups, the upper and lower elementary
unipotents have distinct unique fixed points over the algebraic closure.
Both belong to $\mathrm{PSL}_2(q)$. Hence a pair generating either
projective group cannot share an invariant line. Conjugation by $[X]$
normalizes $\mathrm{PGL}_2(q)$ and its determinant-square subgroup
$\mathrm{PSL}_2(q)$, and its square is the identity.
\end{proof}

For rank-three rotations $a,b$, apply the theorem to the generating pair
$a,ab$. Since $ab$ is an involution, the resulting involutory automorphism
sends $a$ to $a^{-1}$ and fixes $ab$, hence sends $b$ to $a^2b$.
This is exactly the rank-three rotation reversal rule. The retained
rotation-group correspondence therefore excludes a chiral polyhedron
with either projective group as its full automorphism group.
The correspondence is the retained geometric input; the preceding theorem
supplies the complete algebraic reversal. It does not require the product
of the two generators to be an involution or their cyclic subgroups to
have trivial intersection. This recovers Elliott--Leemans--O'Reilly-Regueiro,
Theorem~4.1, by one argument, including the small fields.

For clarity, the source's auxiliary involution count requires a
characteristic qualification. For $u:x\mapsto x+1$, an inverting
projectivity fixes infinity and has the form $x\mapsto ax+b$.
The conjugation equation forces $a=-1$. In odd characteristic all $q$
choices of $b$ give nonidentity involutions. In characteristic two these
are translations; the choice $b=0$ is the identity and the other $q-1$
choices are involutions. Thus the literal count in its Lemma~4.3 is $q-1$
when the characteristic is two. This correction does not affect the
principal nonexistence theorem.

\section{Two affine ends and one field parameter}
Put $L(q)=\mathrm{PSL}_2(q)$ and $P(q)=\mathrm{PGL}_2(q)$.
We retain the standard theorem that their automorphisms are induced by
projectivities and field automorphisms (with the usual small identifications).
No subgroup-classification theorem is needed in the following construction.
\begin{theorem}\label{E.PointedTriples}
Let $F=\mathbb F_q$, let $t\in F\setminus\{0,1,-1\}$ generate $F$ as a
field over its prime field, and put $k=|t|$. Consider the projective matrices
\[
 A=\begin{pmatrix}-t^{-1}&1\\0&1\end{pmatrix},\qquad
 B=\begin{pmatrix}t&0\\0&1\end{pmatrix},\qquad
 C=\begin{pmatrix}1&0\\1+t&-t\end{pmatrix}.
\]
Their pair products $AB,BC$ and triple product $ABC$ are nonidentity
involutions. Set $D=\langle-1,t\rangle\le F^\times$.
The groups $\langle A,B\rangle$ and $\langle B,C\rangle$ are the affine
groups $F^+\rtimes D$ fixing infinity and zero respectively; their
intersection is the diagonal group $D$. Consequently the rank-four
rotation intersection conditions hold exactly when $-1\in\langle t\rangle$.
In that case $\langle A,B,C\rangle$ equals $L(q)$ if $t$ is a square,
and $P(q)$ otherwise; in characteristic two these groups coincide.
The resulting polytope has type
\[
 (|{-t}|,k,|{-t}|).
\]
It is chiral exactly when no field automorphism sends $t$ to $t^{-1}$.
For a fixed admissible $k$, the number of isomorphism classes of the
ordered chiral markings is $\varphi(k)/[F:\mathbb F_p]$.
\end{theorem}
\begin{proof}
The displayed products have trace zero and are nonscalar; direct
multiplication proves the three involution relations in every characteristic.
The first end fixes infinity. It contains $B$ and
$R=AB=\left(\begin{smallmatrix}-1&1\\0&1\end{smallmatrix}\right)$.
The product $RBRB^{-1}$ is translation by $1-t$. Conjugation by $B$
scales translations by $t$, so its translation subgroup contains
$(1-t)\mathbb F_p[t]=F$. Its group of slopes is precisely $D$.
This proves the first affine description. After interchanging zero and
infinity, the second end contains a diagonal of slope $t^{-1}$ and an
affine map of slope $-t$ and translation $1+t$. The same commutator
argument supplies all translations, since $t\ne\pm1$, and proves its
description. An element in both ends fixes zero and infinity, so is
diagonal; both ends contain exactly the slopes $D$ there.

The nonidentity powers of $A$ fix precisely infinity and $t/(1+t)$;
those of $B$ fix infinity and zero. Thus their cyclic intersection is
trivial by the uniqueness of a projectivity on three distinct points.
Likewise $\langle B\rangle\cap\langle C\rangle=1$, since the fixed
points of $C$ are zero and one. Together with the computed end
intersection this proves precisely the rank-four intersection conditions.

The two ends contain the full opposite elementary unipotent groups.
These generate $L(q)$: elementary row operations reduce any determinant-one
matrix to a diagonal matrix, and such a diagonal is also a product of
upper and lower elementary matrices. The determinants of $A,B,C$ have
classes $[-t^{-1}],[t],[-t]$. Under the intersection hypothesis,
$-1\in\langle t\rangle$, their classes generate just $[t]$.
This proves the claimed whole group, without a maximal-subgroup test.
The eigenvalue ratios give the type. In odd characteristic the intersection
hypothesis says $k$ is even; then $|{-t}|=k$ for $4\mid k$ and
$|{-t}|=k/2$ for $k\equiv2\pmod4$. In characteristic two it is $k$.

For chirality use the retained rank-four rotation criterion: a reversing
automorphism would send $A$ to $A^{-1}$, $B$ to $A^2B$ and fix $C$.
Write it as a projectivity after a field automorphism $\sigma$.
Since $C$ is semisimple with eigenvalue ratio $-t$ and has order greater
than two, fixing its projective conjugacy class forces
$t^\sigma\in\{t,t^{-1}\}$. If $t^\sigma=t$, the full-field hypothesis
makes $\sigma$ the identity. A projectivity centralizing $C$ then fixes
zero and one individually, whereas one inverting $A$ must interchange
infinity and $u=t/(1+t)$. There is no such projectivity: writing it as
$\left(\begin{smallmatrix}a&0\\c&d\end{smallmatrix}\right)$, these
conditions give $a=c+d$, $a=uc$ and $d=-uc=-a$. Thus $c=2a$ and
$u=1/2$, forcing $t=1$; in characteristic two they already give $c=0$,
contradicting the interchange. This proves impossibility of inner reversal.

If instead $t^\sigma=t^{-1}$, then $\sigma^2=1$ because $t$ generates
$F$. The matrix $R$ above is defined over the prime field and satisfies
$R^2=I$. Direct substitution gives
\[
 RA^\sigma R=A^{-1},\qquad RB^\sigma R=A^2B,
 \qquad RC^\sigma R=-t^{-1}C.
\]
Thus $X\mapsto RX^\sigma R$ is the required involutory projective
automorphism, and the polytope is directly regular. This establishes
both directions of the chirality assertion.

Finally an isomorphism of two ordered markings sends the first end to
the first end and the second to the second. Their unique fixed points
are infinity and zero, so the projective part of its semilinear action
fixes these individually. It is diagonal, and comparison of $B$ shows
that the new parameter is $t^\sigma$. Conversely field automorphisms
send the whole displayed triple to the triple with parameter $t^\sigma$.
Every such orbit has length $[F:\mathbb F_p]$, and there are
$\varphi(k)$ elements of order $k$. This gives the stated number.
\end{proof}

\begin{proposition}\label{E.PointedCompleteness}
Suppose a rank-four rotation marking generates $L(q)$ or $P(q)$ and each
of its two ends fixes a point of the projective line. Then, up to
projective conjugation, it has the displayed form in
Theorem~\ref{E.PointedTriples}, with $t$ generating the whole field.
\end{proposition}
\begin{proof}
The two fixed points differ, since the whole group fixes no point.
Their common generator $B$ fixes both, so conjugate them to infinity
and zero and write $B=\operatorname{diag}(t,1)$.
The first and last generators have the forms
$A=\left(\begin{smallmatrix}\lambda&\mu\\0&1\end{smallmatrix}\right)$
and $C=\left(\begin{smallmatrix}1&0\\\nu&\kappa\end{smallmatrix}\right)$.
The involution relations give $\lambda=-t^{-1}$, $\kappa=-t$ and
$\mu\nu=1+t$. Here $t\ne1$, since $B$ is nonidentity, and $t\ne-1$:
otherwise $\mu\nu=0$ and the whole group fixes zero or infinity.
Thus $\mu\ne0$, and conjugation by $\operatorname{diag}(1,\mu)$
normalizes $\mu$ to one. This gives the displayed matrices.
They all lie in $P(\mathbb F_p(t))$, so the group-order formulas force
$\mathbb F_p(t)=F$. (For a proper subfield of size $s$, $q\ge s^2$
and $|L(q)|>|P(s)|$.) The preceding theorem applies.
\end{proof}

The full-field condition is $k\nmid p^e-1$ for $0<e<d$, where
$q=p^d$; the absence of field reversal is $k\nmid p^e+1$ for the same
range. Thus this one theorem gives both source pointed-parabolic
classifications, including their counts. For $P(q)$ with odd $q$, add
$k\nmid(q-1)/2$; for $L(q)$ with odd $q$, add $k\mid(q-1)/2$ and
$k$ even. The latter implies $q\equiv1\pmod4$. In characteristic two
there is no determinant-square distinction and no even-$k$ requirement.

There is a proof-route correction to the source's subfield sentence.
When $t$ generates $F$ but a field automorphism inverts it, the group
can still be all of $L(q)$ or $P(q)$. The exclusion from the chiral list
comes from the semilinear reversal just constructed. For example, an
element $t\in\mathbb F_{25}$ of order six gives $L(25)$, of order
$7800$, with affine ends of order $150$ and intersection of order six;
Frobenius inversion makes this marking regular. These orders follow
uniformly from the theorem, not from a sample extrapolation.

\section{Anchored involutions and field descent}
The retained rotation correspondence uses generators
$\sigma_1,\ldots,\sigma_{n-1}$ and the interval-involution relations.
Its intersection axiom says that subgroups generated by interval products
with endpoints in $I$ and $J$ intersect in the subgroup for $I\cap J$.
We use this full axiom, not merely the rank-four special case.

\begin{lemma}\label{E.Anchored}
In the rotation group of a rank-$n$ chiral polytope, the involutions
$t_i=\sigma_i\cdots\sigma_{n-1}$, $1\le i\le n-2$, form a string
C-marking of rank $n-2$. Its adjacent labels are
$|\sigma_1|,\ldots,|\sigma_{n-3}|$, and its even subgroup
$\langle\sigma_1,\ldots,\sigma_{n-3}\rangle$ has index two.
There is a dual anchored marking with the labels at the other end.
\end{lemma}
\begin{proof}
Every $t_i$ is a nonidentity involution by the rotation correspondence.
For $i<j$, cancellation gives $t_it_j=\sigma_i\cdots\sigma_{j-1}$.
If $j-i\ge2$ this is an involution, proving the string commutations;
if $j=i+1$ it is $\sigma_i$. For a set of indices $I$, the group
$\langle t_i:i\in I\rangle$ is exactly the endpoint subgroup of the
rotation intersection axiom with vertex indices
$\{i-1:i\in I\}\cup\{n-1\}$. Products $t_it_j$ supply every
interval with both endpoints in this set. Intersecting two such sets
therefore proves the complete C-intersection property, including the
empty intersection. The even words in the $t_i$ are generated by the
adjacent products $\sigma_1,\ldots,\sigma_{n-3}$. This subgroup cannot
contain $t_{n-2}$: its endpoint set meets that for $t_{n-2}$ in a single
vertex, whose endpoint subgroup is trivial. Thus its index is two.
Reversing the indexing gives the dual statement.
\end{proof}

\begin{lemma}\label{E.NoTwoLabel}
If a chiral rotation group is $L(q)$ or $P(q)$, $q\ge4$, all its labels
are greater than two.
\end{lemma}
\begin{proof}
The first label cannot be two: the identity automorphism would then
satisfy the rotation reversal rule
$\sigma_1\mapsto\sigma_1^{-1}$,
$\sigma_2\mapsto\sigma_1^2\sigma_2$ and fix the other generators.
The same holds at the last label by duality.
Suppose instead $\sigma_j^2=1$ with $1<j<n-1$. Put
$A=\langle\sigma_1,\ldots,\sigma_{j-1}\rangle$ and
$B=\langle\sigma_{j+1},\ldots,\sigma_{n-1}\rangle$.
For $x=\sigma_i\cdots\sigma_{j-1}$ and
$y=\sigma_{j+1}\cdots\sigma_k$, all of $x\sigma_j$, $\sigma_j y$
and $x\sigma_j y$ are involutions. Thus $\sigma_j$ inverts each
such $x,y$, and the last relation gives $xy=yx$. These prefix and
suffix elements generate $A,B$. Hence $A$ and $B$ commute and are
normal in $\langle A,B,\sigma_j\rangle=G$; both are nontrivial.
This is impossible in $L(q)$ or $P(q)$: every nontrivial normal subgroup
contains the simple socle $L(q)$, whose centralizer is trivial.
These are exactly the standard structural inputs already retained in
Topic E. No assumption about a minimal generating number is used.
\end{proof}

\begin{lemma}\label{E.SubfieldEnvelope}
Let $k\subseteq K$ be finite fields of odd characteristic, $|k|\ge5$.
Any conjugate of $P(l)$ or $L(l)$ in $P(K)$ containing the standard
$L(k)$ has $k\subseteq l$ and is the standard $P(l)$ or $L(l)$.
Consequently its intersection with standard $P(k)$ is $P(k)$ in the
first case. In the second case it is $P(k)$ if $[l:k]$ is even and
$L(k)$ if $[l:k]$ is odd.
\end{lemma}
\begin{proof}
The projective invariant $J(X)=\operatorname{tr}(X)^2/\det(X)$ lies
in $l$ on such a conjugate. The matrices
$\left(\begin{smallmatrix}a&-1\\1&0\end{smallmatrix}\right)\in L(k)$
show that $a^2\in l$ for every $a\in k$. Squares span $k$, since
$a=((a+1)/2)^2-((a-1)/2)^2$, so $k\subseteq l$.

Let $\phi$ be the $|l|$-power field automorphism of $K$.
The standard $P(l)$ is its pointwise fixed projective group: normalize
one nonzero matrix entry to one to check this assertion directly.
A conjugate $P(l)^g$ is the fixed group of the conjugate semilinear
automorphism $\theta$. Both $\theta$ and $\phi$ fix $L(k)$ pointwise,
so their quotient is an inner projectivity centralizing $L(k)$.
That centralizer is trivial: it centralizes the upper and lower
nonidentity elementary unipotents, hence fixes infinity and zero,
and commutation with translation forces its diagonal ratio to be one.
Thus $\theta=\phi$ and $P(l)^g=P(l)$. The same argument applied to
an overgroup of $L(l)^g$ places it in standard $P(l)$; its simple
index-two subgroup is standard $L(l)$. Finally a scalar in $k^\times$
is a square in $l$ exactly when it is a square in $k$ or $[l:k]$ is
even. Applying this to determinants proves the intersections.
\end{proof}

We use the accepted Topic E classification: a projective group with a
rank-four C-marking is $P(5)$, $L(11)$ or $L(19)$, with its specified
marking. The proof and finite table belong to the complete common paper
and remain paid once. We also retain Dickson's proper-subgroup alternatives
in their natural-subfield form. Here are the precise additional interfaces.

\begin{lemma}\label{E.RegularSubgroups}
A directly regular rank-three C-group contained in $P(q)$, with both
labels greater than two, is either $S_4$ with its tetrahedral marking,
or a natural conjugate of $P(k)$ with $k$ odd, $|k|\ge5$, whose rotation
subgroup is $L(k)$. A directly regular rank-four C-group contained in
$P(q)$, with all labels greater than two, is the simplex $P(5)\cong S_5$.
No C-marking contained in $P(q)$ with these labels has rank at least five.
\end{lemma}
\begin{proof}
In odd characteristic an affine subgroup has only half-turn involutions;
two distinct such involutions cannot commute. It cannot contain a
rank-three string marking. In characteristic two its involutions are
translations and all commute. A cyclic group is likewise excluded.
In a dihedral group, commuting endpoint involutions of a rank-three
marking cannot be central, so are opposite reflections. Writing them
$s,sr^{N/2}$ and the middle reflection $sr^a$, generation says
$\gcd(a,N/2)=1$. At least one of $a,a-N/2$ is coprime to $N$; the
corresponding adjacent pair already generates the whole group,
contradicting C-intersection. A central middle involution would instead
give a label two.

All involutions in $A_4$ commute. The group $A_5$ has no index-two
subgroup. A directly regular marking of $S_4$ consists of transpositions,
since its even subgroup is its unique index-two subgroup $A_4$.
Commuting distinct endpoints are disjoint transpositions and a middle
transposition joining them gives precisely type $\{3,3\}$.
The remaining Dickson alternatives are natural projective subfield groups.
The simple $L(k)$ has no index-two subgroup; $P(k)$ has one exactly in
odd characteristic, namely $L(k)$. The small fields give the cases already
handled. This proves the rank-three assertion.

For rank four, the affine and dihedral exclusions still apply to the
first three generators. The involution centralizer in $A_5$ is a Klein
four group, so a first generator commuting with the noncommuting last
two is impossible. In $S_4$, if an endpoint were a transposition its
centralizer would be abelian, giving the same contradiction. Thus both
endpoints are double transpositions. The two middle generators must be
transpositions, since adjacent labels exceed two. The first adjacent
pair generates a dihedral group of order eight containing the entire
normal Klein four group, hence the last generator, contrary to the
intersection property. This also excludes $A_4$.
Dickson now reduces to projective subfield groups. The common rank-four
classification leaves $P(5),L(11),L(19)$; only $P(5)$ has an index-two
subgroup, and its marking is the simplex marking.

Finally, at rank at least five both fourth and fifth generators would
centralize the first noncommuting pair. The common projective centralizer
lemma allows at most one nonidentity element there in odd characteristic;
in characteristic two the centralizer of an involution is abelian and
already prevents rank four. This is independent of which subgroup the
marking generates.
\end{proof}

\begin{theorem}\label{E.ChiralRankBound}
Neither $L(q)$ nor $P(q)$, $q\ge4$, is the full group of a chiral
polytope of rank at least five.
\end{theorem}
\begin{proof}
All labels exceed two. Anchoring gives a directly regular C-subgroup of
rank $n-2$, so $n\le6$. If $n=6$, the anchored rank-four subgroup is a
simplex, so the first three labels are three. Its dual makes the last
three labels three. The full rotation group is therefore a quotient of
the even simplex Coxeter group $[3,3,3,3,3]^+\cong A_7$.
This group is simple, so the nontrivial quotient is itself, with its
regular simplex marking. This contradicts chirality.

For $n=5$ write the rotations as $\alpha,\beta,\gamma,\delta$.
The two anchored directly regular groups are
\[
 P_0=\langle\alpha,\beta,\gamma\delta\rangle,
 \qquad Q_0=\langle\alpha\beta,\gamma,\delta\rangle,
\]
with even subgroups $R=\langle\alpha,\beta\rangle$ and
$S=\langle\gamma,\delta\rangle$. If both are tetrahedral,
all four labels are three and the same argument uses the simple
$[3,3,3,3]^+\cong A_6$ and its regular marking.

First suppose the whole group is $P(q)$ with $q$ odd. The even group
of either anchored subgroup lies in $L(q)$: it is a natural $L(k)$,
or is $A_4$ generated by elements of order three, which must lie in
the determinant-square kernel. Thus all four rotations lie in $L(q)$,
a contradiction. If $q$ is even there are no odd-characteristic natural
$P(k)$ subgroups, so both anchored groups would be tetrahedral, already
excluded.

It remains that the whole group is $L(q)$ with $q$ odd, and at least
one anchored group, say $P_0$, is standard $P(k)$, $|k|\ge5$.
Then $R=L(k)$ and $[\mathbb F_q:k]$ is even. Put
$F_0=\langle\alpha,\beta,\gamma\rangle$, a proper subgroup by
rotation intersection. Dickson's alternatives containing $R$ are
projective subfield groups; the sole exceptional possibility is
$R=F_0=A_5$, which would imply
$G=\langle P_0,\gamma\rangle=P_0$, impossible.
The subfield-envelope lemma puts $F_0$ in standard $P(l)$ or $L(l)$
with $k\subseteq l$. The first possibility contains $P_0$, and again
$G=\langle F_0,P_0\rangle=F_0$ is impossible.
In the second, the exact rotation intersection is
$F_0\cap P_0=R=L(k)$. Thus $[l:k]$ is odd, while
$[\mathbb F_q:k]$ is even. It follows that
$[\mathbb F_q:l]$ is even, so $P(l)$ is a proper subgroup of $L(q)$.
But $G=\langle F_0,P_0\rangle\le P(l)$, the final contradiction.
\end{proof}

\section{A trace quadratic for an exceptional end}
Throughout this section the field has odd characteristic. Traces refer to
specified determinant-one lifts. Thus a projective element of order three
has chosen trace $1$, an element of order four has chosen trace
$\sqrt2$, and the two classes of order-five elements have traces
$(-1\pm\sqrt5)/2$. The last convention is used away from characteristic
five. Changing the sign of a lift changes none of the projective statements.

\begin{theorem}\label{E.TraceExtension}
Let $F$ be a finite field of odd order, and let $a,b\in F^\times$, $c\in F$ with
$d=a^2+b^2-4\ne0$. Up to simultaneous projective conjugacy, pairs
$A,B\in\mathrm{SL}_2(F)$ with traces $a,b$ and $\operatorname{tr}(AB)=0$
form one class. For each such pair, determinant-one matrices $C$ with
\[
 \operatorname{tr}(C)=c,\qquad
 \operatorname{tr}(BC)=\operatorname{tr}(ABC)=0
\]
are in bijection with the square roots in $F$ of
\[
 \Delta=a^2c^2-4(a^2+b^2+c^2-4).
\]
There are two, one or zero extensions according as $\Delta$ is a nonzero
square, zero or a nonsquare. Their invariant is
$2\operatorname{tr}(AC)-ac=\pm\sqrt\Delta$.
If $a,c\ne0$ and $c^2\ne4$, a projective conjugation realizes the
rotation reversal precisely when $\Delta=0$.
\end{theorem}
\begin{proof}
Put $J=\left(\begin{smallmatrix}0&1\\-1&0\end{smallmatrix}\right)$.
Choose $U,V\in F$ with $U^2+V^2=d$. Such a choice exists because the
sets of squares and of $d$ minus squares each have $(|F|+1)/2$ elements.
Set $u=(U+a)/2$, $w=(V-b)/2$ and
\[
 A=\begin{pmatrix}u&w+b\\w&a-u\end{pmatrix},\qquad B=A^{-1}J.
\]
Direct expansion gives determinant one, the required traces, and $AB=J$.
The nonzero $d$ says that $A,B$ have no common eigenline: equivalently,
$\det(AB-BA)=-d$. More generally the four matrices $I,A,B,AB$ are
linearly independent. This follows by diagonalizing a nonscalar $A$:
both off-diagonal entries of $B$ are nonzero, so $I,A$ span the diagonal
matrices and $B,AB$ supply the off-diagonal ones. In the remaining Jordan
case its unique eigenline is not $B$-invariant, and the same four-entry
calculation proves independence.
The Cayley--Hamilton identities and
$AB+BA=bA+aB-abI$ determine their multiplication table from $a,b$.
Consequently any second pair with the same three traces defines an
algebra automorphism of $M_2(F)$ by matching these bases. Such an
automorphism is inner: take a nonzero vector in the image of $E_{11}$,
and its image under the transformed $E_{21}$, as a basis. The transformed
matrix units then act as the usual matrix units. This proves pair uniqueness.

Write $C=\left(\begin{smallmatrix}r&s\\s&c-r\end{smallmatrix}\right)$,
as is forced by $\operatorname{tr}(JC)=0$, and put $R=2r-c$, $S=2s$.
The other trace and determinant equations are exactly
\[
 VR-US=-bc,\qquad R^2+S^2=c^2-4.
\]
Since $d\ne0$, every solution is uniquely of the form
\[
 R=\frac{-bcV+zU}{d},\qquad S=\frac{bcU+zV}{d},\qquad z^2=\Delta.
\]
This also gives $2\operatorname{tr}(AC)-ac=UR+VS=z$.
The projective centralizer of the ordered pair is trivial, so the
extensions have not acquired an additional conjugacy identification.

For the last assertion, nonzero traces remove the signs in determinant-one
lifts of $XAX^{-1}=A^{-1}$ and $XCX^{-1}=C$.
Put $A_0=A-aI/2$, $C_0=C-cI/2$. The centralizer algebra of $C$ is
$F[I,C_0]$, so $X=xI+yC_0$. The first equation is
\[
 XA_0+A_0X=2xA_0+y\operatorname{tr}(A_0C_0)I=0.
\]
Since $A_0$ is nonscalar, invertibility forces $x=0$, $y\ne0$ and
$\operatorname{tr}(A_0C_0)=z/2=0$. Conversely, if $z=0$, choose $X=C_0$.
It is invertible because $c^2\ne4$, has scalar square, commutes with $C$,
and inverts $A$ projectively. Also $C_0J=-JC_0$ by the trace condition,
so it fixes $AB=J$ projectively and sends $B$ to $A^2B$.
This is exactly the rotation reversal rule.
\end{proof}

\begin{lemma}\label{E.OddDegreeExclusion}
If $q=p^d\equiv3\pmod4$ and $d>1$ is a prime power, then $L(q)$
cannot be a chiral rotation group.
\end{lemma}
\begin{proof}
Only rank four remains. Neither end group fixes a point: it contains the
involution given by the product of its rotations, whereas an involution
of $L(q)$ has no rational fixed point when $q\equiv3\pmod4$.
By the retained natural-subfield maximal-subgroup classification, each
proper end lies in a conjugate of $L(p^{e_i})$ with $e_i\mid d$ proper.
Here the exceptional $A_4,S_4,A_5$ alternatives also lie in proper
subfield overgroups; the odd degree excludes a natural $P(p^{e_i})$.
Conjugate the first end to its standard position and choose lifts $A,B$
there. A lift of the third rotation has trace $c\in F_{p^{e_2}}$, since
trace is conjugacy invariant. Put $F=F_{p^{\operatorname{lcm}(e_1,e_2)}}$;
this is a proper subfield because the divisors of the prime power $d$
are linearly ordered.

The first pair is absolutely irreducible. Otherwise its common
algebraic eigenline is rational, contrary to the preceding fixed-point
observation, or its Frobenius conjugate is a distinct common eigenline.
In the latter case the pair is diagonal on the two lines and commutes.
Then $(AB)^2=1$ makes $A^2=B^{-2}$ projectively; the cyclic intersection
axiom forces both rotations to have order at most two, already excluded.
Thus $a^2+b^2-4\ne0$. Normalize $AB$ to $J$ over $F$ and apply the
explicit extension formulas. All entries of $C$ lie in $F(\sqrt\Delta)$.
The ambient extension has odd degree, so this quadratic extension is
$F$ itself. Hence all three rotations lie in the proper $L(F)$,
a contradiction. This replaces the source's quadratic calculation,
including its printed coefficient errors, by a nonsingular formula.
\end{proof}

\section{Exceptional ends: three discriminants and the exact intersections}
The finite input to this section is fully specified, rather than hidden
behind a group-name lookup. The program \texttt{exceptional\_amalgams.g}
uses exact GAP subgroup and finitely presented group algorithms. It first
enumerates every ordered generating pair $(a,b)$ in $A_4,S_4,A_5$ with
$|a|,|b|>2$, $|ab|=2$ and $\langle a\rangle\cap\langle b\rangle=1$.
The numbers of pairs are respectively $24,48,360$, and their types modulo
automorphisms are
\[
 A_4:(3,3),\qquad S_4:(3,4),(4,3),\qquad
 A_5:(3,5),(5,3),(5,5).
\]
For every $(5,5)$ pair in $A_5$ it checks
$a b^2a^3b^2a^2=1$. This is a complete finite premise, not a sample of
projective groups. The standard spherical triangle presentations identify
the $(3,3),(3,4),(3,5)$ groups as $A_4,S_4,A_5$.

\begin{lemma}\label{E.ExceptionalTypes}
If $p\equiv3\pmod4$ is prime and $L(p)$ is a chiral rank-four rotation
group, its type, up to duality, is $(5,3,4)$, $(5,3,5)$ or $(3,5,3)$.
\end{lemma}
\begin{proof}
Neither proper end fixes a point. A dihedral end with both rotation
orders greater than two is impossible: both rotations would lie in its
cyclic rotation subgroup, and their trivial cyclic intersection together
with $(ab)^2=1$ would force their orders to be at most two. The retained
prime-field subgroup theorem therefore leaves $A_4,S_4,A_5$.
The displayed complete local pair census leaves only the compatible
three-label types. The types $(3,3,3)$, $(3,3,4)$, $(3,4,3)$ and
$(3,3,5)$ and their duals are excluded by their spherical universal
rotation groups. These have respectively order $60,192,576,7200$;
the latter has only $C_2$ and $A_5$ composition factors, and the middle
two are solvable. No simple quotient is $L(p)$ with $p\equiv3\pmod4$.
These are the standard simplex, cube, $24$-cell and $120$-cell groups;
their full presentation orders and composition factors are also checked
by the same finite program.

For $(4,3,4)$ the trace quadratic has $a^2=c^2=2$, $b^2=1$ and
$\Delta=0$. Its explicit projective reversal excludes chirality.
The remaining unwanted types are $(5,5,3)$ and $(5,5,5)$ and their duals.
Add the preceding five-five relator to each five-five end of its full
rotation presentation. The exact finite calculation gives order $7200$
and a composition series with successive orders $7200,120,2,1$ in both
cases. Their only nonabelian simple composition factor is $A_5$.
Again $L(p)$ cannot be a quotient. This exhausts every compatible local
pair, without an unrecorded search in $L(p)$.
\end{proof}

For completeness we give the full finite intersection test. Write
\[
 T(r,s,t)=\langle a,b,c\mid a^r=b^s=c^t=(ab)^2=(bc)^2=(abc)^2=1\rangle.
\]
In an $A_5$ end of type $(5,3)$ or $(3,5)$ the word
$u(a,b)=ab^2a^2b^2a^2$ is an involution inverting $b$ and generates,
with $b$, its full dihedral normalizer. In an $S_4$ end of type $(3,4)$
the corresponding word is $v(c,b)=cb^2c^2$. All these word assertions
are checked on every relevant local pair. The only proper overgroups
of $\langle b\rangle$ in $A_5$ have orders $3,6,12$ when $|b|=3$,
and $5,10$ when $|b|=5$. Thus a dihedral collision between the two ends
forces $u(c,b)=u(a,b)b^k$ (using $v$ in the $S_4$ case). For
$k=1,\ldots,s$ the exact quotient orders are
\[
\begin{array}{c|c}
 (r,s,t)& |T(r,s,t)/\langle\!\langle\hbox{collision}\rangle\!\rangle|\\\hline
 (5,3,4)&60,1,1\\
 (5,3,5)&3420,60,60\\
 (3,5,3)&660,1,360,360,1.
\end{array}
\]
For an $A_4$ collision put
$u_1=a^2b^{-1}a$, $u_2=ab^{-1}a^2$.
The two subgroups $\langle u_i,b\rangle$ are exactly the $A_4$
overgroups in the first $A_5$. Their three involutions are the
$b$-conjugates of $u_ib$. At the other end use $v=b^2c^2$ for $S_4$,
or $v_1=c^2b^{-1}c$, $v_2=cb^{-1}c^2$ for $A_5$.
It is therefore enough, and is exhaustive, to add each relation
$u_i^{b^j}=v_l$, $j=0,1,2$.
The six $(5,3,4)$ quotients all have order two.
For $(5,3,5)$ the twelve orders, in lexicographic $(i,l,j)$ order, are
\[
 1,1,1,\quad1,60,60,\quad60,1,60,\quad1,1,1.
\]
This corrects the source's assertion in Lemma~5.12 that all twelve
quotients are trivial. Its exclusion of $A_4$ intersections outside
$L(4)\cong L(5)$ remains valid. Full finitely presented group orders,
not orders of selected permutation images, are used in this calculation.
The nontrivial collision groups are identified as $A_5$, $L(9)$, $L(11)$
and $L(19)$ at orders $60,360,660,3420$. Their simplicity shows that the
complete intersection conclusions hold for arbitrary projective quotients:
exclude $q=4,5$ for $(5,3,4)$; $q=4,5,19$ for $(5,3,5)$; and $q=9,11$
for $(3,5,3)$. Thus the source's three intersection lemmas, not merely
their prime-field applications, follow from this one finite interface.

\begin{theorem}\label{E.PrimeExceptional}
Let $p\ge7$ be prime. For each square root of $5$ in $F_p$ define
\[
 D_{534}=1\pm\sqrt5,\qquad
 D_{535}=(7\pm5\sqrt5)/2,\qquad
 D_{353}=3\pm2\sqrt5.
\]
For type $(5,3,4)$ also require that $2$ is a square in $F_p$.
For any of these three types, the number of ordered chiral polytope
isomorphism classes with rotation group $L(p)$ is twice the number
of nonzero squares in the corresponding two-element list.
In the $(3,5,3)$ case this formula gives zero at $p=11$; in the
$(5,3,5)$ case it gives two at $p=19$. A missing square root of $5$
means there are no such classes.
The same types with group $L(q)$, $q\equiv3\pmod4$, require $q=p$ prime.
\end{theorem}
\begin{proof}
A local $(5,3)$ or $(3,5)$ pair generates $A_5$ by the spherical
triangle presentation and simplicity. A $(3,4)$ pair generates $S_4$:
the exact order four prevents its proper triangle quotients. Fix
traces as above. The nonzero quantity $d=a^2+b^2-4$ follows either
from the absolute irreducibility of the $A_5$ end or directly from
the quadratic equation for its order-five trace. The trace extension
theorem gives all extensions and their conjugacy classes.
For $(5,3,4)$ it gives $D_{534}$, while for $(3,5,3)$ it gives $D_{353}$.
For $(5,3,5)$, if $a^2=c^2$, it gives $D_{535}$.
If the two order-five trace squares differ, their sum is three and
product one, so $\Delta=1$. Both resulting extensions lie inside the
first $A_5$. Indeed, for the fixed pair
$a=(1\,2\,3\,4\,5)$, $b=(2\,5\,3)$ the complete local enumeration
finds exactly two order-five elements $c$ satisfying both involution
relations, and both are in the other $A_5$ conjugacy class from $a$.
The two classes have different trace squares in every characteristic
under consideration: replacing an order-five element by its square
replaces its trace square $t$ by $(t-2)^2=3-t$, using
$t^2-3t+1=0$; these differ because the characteristic is not five. Pair uniqueness and the two-root formula show
that these are exactly the mixed-class extensions. They cannot generate
$L(p)$.

Every remaining extension generates $L(p)$. A subgroup containing
$A_5$ in a prime-field $L(p)$ is either that $A_5$ or the whole group,
by the retained subgroup classification. In type $(5,3,4)$, $c$ has
order four and hence is outside $A_5$. In type $(5,3,5)$ the preceding
complete local check excludes same-class $c$ inside $A_5$.
In type $(3,5,3)$ the complete local check finds no order-three $c$
inside the first $A_5$ satisfying the two involution relations.
Thus generation is uniform; no large-group enumeration is substituted
for it.

The intersection table proves the end intersection is exactly
$\langle b\rangle$, except possibly at $p=19$ in type $(5,3,5)$
and $p=11$ in type $(3,5,3)$. Adjacent cyclic intersections are trivial
by the coprime orders. The required program \texttt{exceptional\_small.g}
checks every root of the same equations at these two primes. At $p=19$
the two nonzero-root extensions have end intersection of order three;
the zero-root extension has intersection six. At $p=11$ the sole
$(3,5,3)$ extension has zero discriminant and intersection ten.
Thus every nonzero-root extension has the full rank-four rotation
intersection property, and none of the zero-root cases contributes.

For a prime field every automorphism of $L(p)$ is induced by a
projectivity, a retained classical input. The trace extension theorem
therefore rules out reversal for every nonzero root. It produces the
reversal at a zero root whether or not the intersection property holds.
The two choices of order-five trace square cannot be identified by
an automorphism, and the invariant $2\operatorname{tr}(AC)-ac$
distinguishes the two roots within a choice. Conversely pair uniqueness
shows that these are all classes. This proves the exact counts.

Finally suppose $q=p^d\equiv3\pmod4$ and such a type is realized.
Here $d$ is odd. Its local $A_5$ (and $S_4$ where applicable) traces
lie in an extension of the prime field of degree at most two, so they
already lie in $F_p$. The same applies to a square root of $\Delta$
inside $F_q$. Conjugate the first pair to the displayed prime-field
normal form. The unique extension formulas put the third matrix in
$\mathrm{SL}_2(p)$ as well. Hence the generated group is conjugate to
$L(p)$, and equality with $L(q)$ forces $d=1$.
\end{proof}

The residue conditions of the source now follow from ordinary quadratic
reciprocity. The products of the paired discriminants are respectively
$-4,-19,-11$. For $p\ne19,11$ the Legendre symbols of the latter two
are $(p/19)$ and $(p/11)$. The quadratic residues modulo $19$ are
$1,4,5,6,7,9,11,16,17$; those modulo $11$ are $1,3,4,5,9$.
If the product is a nonsquare exactly one discriminant is a square;
if it is a square both are squares or neither is.
For $(5,3,4)$ the order-five and order-four trace requirements give
$p\equiv1,9,31,39\pmod{40}$. At residues $31,39$ exactly one
$D_{534}$ is a square, and at $1,9$ either both or neither are.
These are the complete counts of source Theorems~6, 7, 9 and~11.
In particular residue $10\pmod{11}$ belongs to the nonsquare-product
case of Theorem~11. Its omission from Theorem~1(4b) in the exact
arXiv-v1 edition is a statement error, not an additional exception to
the construction. A separately recorded $p=571$ matrix witness makes
that omission explicit.

\begin{proposition}\label{E.AllPointedExistence}
For $q\ge4$, $P(q)$ occurs as the full group of a chiral rank-four
polytope exactly when $q>4$. For $q\equiv1\pmod4$, $L(q)$ occurs
exactly when $q>9$.
\end{proposition}
\begin{proof}
For $P(q)$ choose a primitive $t$, of order $q-1$. If $q>4$, its order
exceeds $p^{d-1}+1$ when $d>1$, so no nonidentity field automorphism
fixes or inverts it. The pointed theorem applies, with nonsquare $t$
in odd characteristic. For $L(q)$, $q\equiv1\pmod4$, choose $t$ of
order $(q-1)/2$. This is even and greater than two when $q>9$.
For $d>1$ it exceeds $p^{d-1}+1$ except at $p^d=9$: the required
strict inequality is $p^{d-1}(p-2)>3$. Again the pointed theorem
applies. At $d=1$ there is no field automorphism to test.

The omitted groups are $L(4)=P(4)=L(5)\cong A_5$ and $L(9)\cong A_6$.
The complete finite proof in \texttt{small\_rotation\_exclusions.g}
examines all pairs $(a,b)$ modulo the actual automorphism group, retaining
$|a|,|b|>2$, $|ab|=2$, trivial cyclic intersection and a proper first end.
For each it examines every $c=(ab)t$ with $t$ an involution; this is
exactly the condition $(abc)^2=1$. It then applies the remaining
relations, generation and complete rank-four intersection conditions.
There are one and seven pair representatives, respectively. The resulting
polytopal triples number one and four (before the final stabilizer
quotient), and each has a verified involutory automorphism
$(a,b,c)\mapsto(a^{-1},a^2b,c)$. Thus none is chiral. This is a complete
finite premise for these groups, separate from the uniform existence proof.
\end{proof}

\section{Primitive squares in two odd extensions}
This section reproduces the project's earlier accepted primitive-square
proof, independently audited in phase267. Its new placement in the
common theory is not a new answer to a source question. The retained
analytic input is the multiplicative-character bound in
Cohen--Kapetanakis--Reis, \emph{The existence of $F_q$-primitive points
on curves using freeness}, Theorem~5 (printed p.~643), together with
the cyclic-group primitive-element indicator. The exact primary and
prior proof bindings accompany the scope map. The following particular
norm construction and finite reduction are the previously proved work.

\begin{theorem}\label{E.PrimitiveSquare}
Let $p$ be an odd prime and $e_1,e_2>1$ odd integers. There are primitive
$j_i\in F_{p^{e_i}}$ such that, with $u=j_1+j_1^{-1}$ and
$v=j_2+j_2^{-1}$,
\[
 u^2v^2-4(u^2+v^2)
\]
is a nonzero square in $F_{p^{\operatorname{lcm}(e_1,e_2)}}$.
The degrees may be equal or nested.
\end{theorem}
\begin{proof}
Reorder the degrees and put $g=\gcd(e_1,e_2)$, $q_0=p^g$,
$a=e_1/g\le b=e_2/g$. Then $a,b$ are coprime and odd. Write
$E=F_{q_0^a}$, $K=F_{q_0^b}$, $L=F_{q_0^{ab}}$.
Fix any primitive $j_1\in E$ and set $s=(j_1+j_1^{-1})^2$.
One has $s\ne0,4$, since otherwise $j_1$ has order at most four.
Also $s$ has degree $a$ over $F_{q_0}$: if it lies in a proper
intermediate field, the equations $u^2=s$ and $j_1^2-uj_1+1=0$
put first $u$ and then $j_1$ in that same field, because the remaining
extension degree is odd. This contradicts primitiveness.

Define
\[
 f_s(X)=(s-4)(X^2+1)^2-4sX^2,
 \qquad P(X)=\prod_{i=0}^{a-1}f_{s^{q_0^i}}(X).
\]
For $x\ne0$ the desired quantity with $j_2=x$ is $f_s(x)/x^2$.
The quadratic in $X^2$ defining $f_s$ has nonzero constant and leading
coefficients $s-4$ and discriminant $64s\ne0$. Thus $f_s$ has four
distinct nonzero roots. Distinct quartics have disjoint roots: subtracting
gives $(s-t)(X^2-1)^2$, while $f_s(\pm1)=-16\ne0$.
Consequently $P\in F_{q_0}[X]$ is squarefree of degree $D=4a$ and
$P(0)\ne0$. Since $a,b$ are coprime, the $q_0^b$-Frobenius permutes
all these $a$ conjugates. For $x\in K$ this proves
\[
 N_{L/K}(f_s(x))=P(x),\qquad
 \chi_L(f_s(x)/x^2)=\chi_K(P(x)),
\]
where $\chi$ is the quadratic character extended by zero.
If $a>1$, $P$ has no root in $K$: a root would give
$s^{q_0^i}=4(x^2+1)^2/(x^2-1)^2\in E\cap K=F_{q_0}$,
contradicting its degree. When $a=1$, at most four roots remain.

Put $Q=|K|$, $n=Q-1$, $W(n)=2^{\omega(n)}$ and
$\vartheta(n)=\varphi(n)/n$. For every multiplicative character $\psi$,
\[
 \left|\sum_{x\in K^\times}\psi(x)\chi_K(P(x))\right|
 \le D\sqrt Q.                                             \tag{20}
\]
Indeed, for nontrivial $\psi=\eta^k$, $1\le k<n$, with $\eta$ a
generator of the character group, the summand is
$\eta(x^kP(x)^{n/2})$. Its polynomial has $D+1$ distinct roots and
at a root of $P$ has multiplicity exactly $n/2$, so is not a scalar
times an $n$th power. The retained bound gives $D\sqrt Q$.
For trivial $\psi$, apply the quadratic-character bound on all of $K$
and remove zero, giving $(D-1)\sqrt Q+1\le D\sqrt Q$.
All characters, including the trivial one, are extended by zero at zero.

The primitive-element indicator is
\[
 {\bf1}_{\rm primitive}(x)=\vartheta(n)
 \sum_{d\mid n}\frac{\mu(d)}{\varphi(d)}
       \sum_{\operatorname{ord}\psi=d}\psi(x).
\]
Character orthogonality in the cyclic group $K^\times$ proves this
formula. There are $\varphi(d)$ characters of order $d$ and $W(n)$
squarefree divisors. If $Z$ counts primitive roots of $P$, (20) therefore
makes the number of primitive $x$ with $P(x)$ a nonzero square positive
provided
\[
 Q-1>D W(Q-1)\sqrt Q+Z/\vartheta(Q-1).
\]
For $a>1$ it suffices that
\[
 Q-1>4aW(Q-1)\sqrt Q;                                    \tag{21}
\]
for $a=1$, since $Z\le4$ and $1/\vartheta(n)\le W(n)$, it suffices that
\[
 Q-1>4W(Q-1)(\sqrt Q+1).                                 \tag{22}
\]

We give the full finite reduction. For every positive integer $n$,
\[
 W(n)\le \frac{64}{30030^{1/4}}n^{1/4}<5n^{1/4}.
\]
Each prime at least $17$ contributes $2\le\ell^{1/4}$; only
$2,3,5,7,11,13$, whose product is $30030$, can enlarge the ratio.
Repeated prime powers decrease it, and $64^4<5^4\cdot30030$.
Thus (21) holds when $Q>(21a)^4$.
In the nonnested case $3\le a\le b-2$, $b\ge5$. For odd $b\ge23$,
$3^b>[21(b-2)]^4$, first by direct comparison at $23$ and then because
the ratio increases on replacing $b$ by $b+2$ by a factor at least
$9(21/23)^4>1$. For odd $5\le b\le21$, it remains only to examine
$q_0^b\le[21(b-2)]^4$ with $q_0=p^g$, $g$ odd.
The possible bases are at most $27,13,9,5,5,3,3,3,3$ in this order;
the only composite base of odd prime-power exponent is $27$.
The exact integer comparison
\[
 (Q-1)^2>16a^2W(Q-1)^2Q
\]
for all $66$ remaining parameter triples leaves precisely
\[
 (q_0,a,b)=(3,3,5),(5,3,5),(3,3,7),(3,5,7),(3,5,9),(3,7,9).
\]
This exhaustive bounded arithmetic is part of the required program,
using exact trial division and integer comparisons.

For $a=1$, (22) follows from $Q>21^4$: writing $t=Q^{1/4}>21$,
$(t-20)t^3>t^3>20t+1$ gives $Q-1>20t^3+20t$.
There are exactly $23$ odd prime-power fields $p^e\le21^4$, $e\ge3$
odd. The exact tests $Q-1-4W>0$ and
$(Q-1-4W)^2>16W^2Q$ leave only $Q=3^3,3^5,5^3,7^3,11^3$.
Their exponents are prime, so the divisor $e_1>1$ of $e_2$ equals $e_2$.

Finally put $z_e=Z(p^e)$ in the bound Conway-field representation,
choose $j_1=z_{e_1}$ and $j_2=z_{e_2}^k$, and use the following complete
exceptional list. The displayed exponent is even and nonzero modulo
$p^{e_2}-1$, and $k$ is coprime to that integer.
{\renewcommand{\arraystretch}{1.4}
\[
\begin{array}{ccc|r|r}
 p&e_1&e_2&k& P(j_2)\text{ or }f_s(j_2)\\\hline
 3&3&5&7&z_5^{228}\\5&3&5&1&z_5^{2802}\\
 3&3&7&1&z_7^{1336}\\3&5&7&1&z_7^{406}\\
 3&5&9&5&z_9^{18778}\\3&7&9&1&z_9^{17426}\\
 3&3&3&1&z_3^6\\3&5&5&3&z_5^{200}\\
 5&3&3&1&z_3^{84}\\7&3&3&1&z_3^6\\11&3&3&3&z_3^{906}
\end{array}
\]}
The complete polynomial coefficients and defining irreducible polynomials
are carried in \path{psl2_square_witnesses.json}. The required
\path{verify_square_witnesses.py} checks the finite-field arithmetic
in explicit polynomial quotient rings, independently of GAP; the
required \path{psl2_square_bound_check.py} checks both bounded
reductions. The old GAP witness generators are retained separately as
validation. These eleven exact certificates close all exceptions to
the uniform inequality and hence all parameters of the theorem.
\end{proof}

\section{Separating two subfield ends}
The following is the complete previously proved intersection argument,
now placed next to the common trace obstruction it uses. Its earlier
independent acceptance and finite diagnostic are provenance, not an
unpriced substitute for the proof below.

\begin{lemma}\label{E.TranslatedSubfields}
Let $F_1,F_2$ be subfields of a field $K$, with $F_1\cap F_2=F$.
If $t\in K\setminus(F_1+F_2)$, then
\[
 L(F_1)\cap U(t)L(F_2)U(-t)
 =\{[U(b)]:b\in F\},\qquad U(t)=\begin{pmatrix}1&t\\0&1\end{pmatrix},
\]
in every characteristic.
\end{lemma}
\begin{proof}
Write a determinant-one representative in the intersection as
$M=\left(\begin{smallmatrix}a&b\\c&d\end{smallmatrix}\right)$ over $F_1$.
Its representative over the other field differs only by a scalar $\pm1$,
which can be absorbed. Thus
\[
 U(-t)MU(t)=\begin{pmatrix}
 a-tc&b+t(a-d)-t^2c\\c&d+tc
 \end{pmatrix}
\]
has entries in $F_2$. Hence $c\in F$. If $c\ne0$, the first entry
would put $t$ in $F_1+F_2$, so $c=0$. Now $a,d\in F$; if $a\ne d$,
the upper-right entry gives the same contradiction. Thus $a=d=\pm1$
and $b\in F$. Conversely every displayed unipotent belongs to both
groups and commutes with $U(t)$.
\end{proof}

\begin{theorem}\label{E.CoprimeConstruction}
Let $p$ be odd and $e_1,e_2>1$ coprime odd integers. Put
$F_i=F_{p^{e_i}}$, $K=F_{p^{e_1e_2}}$. Choose primitive $j_i\in F_i$
by Theorem~\ref{E.PrimitiveSquare}, put $u=j_1+j_1^{-1}$,
$v=j_2+j_2^{-1}$ and choose $D\ne0$ with
$D^2=u^2v^2-4(u^2+v^2)$. The projective matrices
\[
 \alpha=\frac12\begin{pmatrix}u&u+2\\u-2&u\end{pmatrix},\quad
 \beta=\frac12\begin{pmatrix}2+u&u\\-u&2-u\end{pmatrix},\quad
 \gamma=\frac1{2u}\begin{pmatrix}v(u-2)&D\\D&v(u+2)\end{pmatrix}
\]
give a chiral rank-four polytope with full group $L(K)$ and type
$((p^{e_1}-1)/2,p,(p^{e_2}-1)/2)$.
\end{theorem}
\begin{proof}
The odd-degree argument used above shows that $u^2$ and $v^2$ generate
their full respective fields over $F_p$; in particular $u,v\ne0$ and
$u^2,v^2\ne4$. The matrices have determinant one, with traces $u,2,v$.
Their eigenvalues give the first and last orders, while $\beta=I+N$
with $N\ne0$, $N^2=0$, gives order $p$. These orders imply the adjacent
cyclic intersections are trivial. Direct multiplication gives
$\alpha\beta=J$ and zero traces for $\beta\gamma,\alpha\beta\gamma$,
so all interval-involution relations hold.

We prove the remaining intersection without assuming generation.
If $x\in F_1$, $y\in F_2$ and $xy\in F_1+F_2$, then $x\in F_p$
or $y\in F_p$. Indeed, apply the $p^{e_1}$-Frobenius $\tau$ to
$xy=a+b$, $a\in F_1,b\in F_2$, and subtract. If $y^\tau\ne y$,
then $x=(b^\tau-b)/(y^\tau-y)\in F_2$, hence in $F_p$; otherwise
$y\in F_p$. Put $z=D/(uv)$. It satisfies
$z^2=1-4/u^2-4/v^2$ and does not belong to $F_1+F_2$.
For if $z=x+y$ with $x\in F_1,y\in F_2$, the square equation puts
$2xy$ in that sum. The preceding observation forces $z$ into one of
the two fields, and then puts $u^2$ or $v^2$ in $F_p$, a contradiction.

With $M^T=T^{-1}MT$, put
\[
 T=\begin{pmatrix}1&2/u\\-1&0\end{pmatrix},\qquad
 A(w)=\begin{pmatrix}1&(2-w)/w\\-w&w-1\end{pmatrix},\qquad
 t=\frac D{2uv}+\frac1v-\frac12-\frac1u.
\]
Direct multiplication gives
\[
 \beta^T=U(1),\quad\alpha^T=A(u),\qquad(\gamma^T)^{U(t)}=A(v),
\]\[
 \gamma^T=\begin{pmatrix}
 (-D+v(u+2))/(2u)&-D/u^2\\-v&(D+v(u-2))/(2u)
 \end{pmatrix}.
\]
The last identity also follows from its upper-left entry $1$, lower-left
entry $-v$, trace $v$ and determinant one. Since $t\notin F_1+F_2$,
Lemma~\ref{E.TranslatedSubfields} puts the two actual end groups'
intersection inside $\langle U(1)\rangle$. Both contain that group,
so their intersection is exactly $\langle\beta\rangle$.

For generation use the retained maximal-subgroup list. The unipotent
$\beta^T$ fixes only infinity, whereas $A(u)$ does not fix infinity,
so the group is in no point stabilizer. The dihedral maximal orders
$|K|\pm1$ are prime to $p$ and cannot contain $\beta$. The exceptional
$A_4,S_4,A_5$ have no element of order at least $13$, whereas
$|\alpha|=(p^{e_1}-1)/2\ge13$. Finally the conjugacy invariant
$\operatorname{tr}(M)^2/\det(M)$ takes the values $u^2,v^2$ on
$\alpha,\gamma$. A conjugate subfield group containing them must
therefore contain both $F_1,F_2$ in its field of definition, and hence
all of $K$. No proper maximal subgroup remains.

A reversing automorphism is semilinear by the retained automorphism
theorem. Its field part fixes $u^2$ and $v^2$, so fixes $K$ and is
trivial. The trace obstruction in Theorem~\ref{E.TraceExtension}, or
its final two-line centralizer calculation alone, now applies:
\[
 \operatorname{tr}(\alpha\gamma)-uv/2=D/2\ne0.
\]
Thus no projectivity can simultaneously invert $\alpha$ and fix
$\gamma$. There is no reversal, so the full rotation correspondence
gives the asserted chiral polytope.
\end{proof}

Every odd integer $d>1$ with at least two distinct prime divisors admits
a factorization $d=e_1e_2$ into coprime odd integers greater than one,
by separating its prime-power factors. This theorem supplies all such
exponents. Together with the earlier exclusions, the nonprime-field
classification is consequently: $L(p^d)$, $d>1$, occurs in rank four
exactly for even $q>4$, for $q\equiv1\pmod4$ with $q>9$, or for
$q\equiv3\pmod4$ when $d$ is not a prime power. This reproduces the
previously accepted Section~6 completion. It is kept separate from
the source's established principal theorem and earns no new question
or historical-priority claim in this edition.


\section{Attribution and computational boundary}
The original three applications retain Leemans--Schulte's two projective
polytope classifications and Leemans--Toledo's maniplex results. Their
complete local and universal table, geometry and literal correction
witnesses remain in \texttt{required\_finite.g}. The rotation application
recovers Elliott--Leemans--O'Reilly-Regueiro and Leemans--Moerenhout--
O'Reilly-Regueiro, with edition-specific corrections made explicit.
The primitive-square and coprime-subfield completion reproduces the
project's earlier accepted work; it is not new question credit.

The new required programs are the complete exceptional amalgam and small
prime certificates, the exhaustive $A_5/A_6$ exclusions, and the unchanged
primitive-square integer reduction and polynomial witness verifier with
its full JSON data. All whole files accompany the common paper. Finite
samples of the pointed construction, the explicit $571$ source correction,
the earlier GAP witness generators and the subfield-intersection example
are validation, with their whole files disclosed separately. No sampled
large projective group supplies a uniform theorem.

Nine new Lean declarations prove the actual $2\times2$ intertwining
identities and the bijective algebra behind the trace quadratic. Existing
common normalizer, matrix and extension certificates, six primitive-square
polynomial declarations and the actual subfield-entry statements are
retained with their exact scopes. Irreducibility and finite-field existence,
Dickson's theorem, projective automorphisms, the character bound, GAP's
finite-presentation computations and the geometric correspondence remain
written or retained inputs. No full Lean proof of a source is asserted.
The source map distinguishes the principal contracts from quoted all-field
Coxeter-quotient classifications and motivating atlas tables not regenerated
here. Full external proof and catalogue costs are unknown, not zero.
\begin{thebibliography}{9}
\bibitem{LS1} D. Leemans and E. Schulte, \emph{Groups of type $L_2(q)$ acting
on polytopes}, arXiv:math/0606660v1; Advances in Geometry 7 (2007), 529--539.
\bibitem{LS2} D. Leemans and E. Schulte, \emph{Polytopes with groups of type
$\PGL_2(q)$}, arXiv:0909.1991v1.
\bibitem{LT} D. Leemans and M. Toledo, \emph{Maniplexes with automorphism
group $\PSL(2,q)$}, eleven-page author manuscript, 15 May 2023.
\bibitem{ELO} D. Elliott, D. Leemans and E. O'Reilly-Regueiro,
\emph{Chiral polyhedra and projective lines}, Int. J. Math. Stat.
16(1) (2015), 45--52; bound eight-page authored edition.
\bibitem{LMO} D. Leemans, J. Moerenhout and E. O'Reilly-Regueiro,
\emph{Projective linear groups as automorphism groups of chiral polytopes},
arXiv:1606.08017v1, 26 June 2016; exact twenty-five-page edition.
\bibitem{SW} E. Schulte and A. I. Weiss, \emph{Chiral polytopes},
Applied Geometry and Discrete Mathematics, DIMACS 4 (1991), 435--516;
rotation-group correspondence.
\bibitem{MS} P. McMullen and E. Schulte, \emph{Abstract Regular Polytopes},
Cambridge University Press, 2002; string C-group correspondence.
\bibitem{CKR} S. D. Cohen, G. Kapetanakis and L. Reis,
\emph{The existence of $F_q$-primitive points on curves using freeness},
C. R. Math. 360 (2022), 641--652, Theorem 5 and formula (2), p.~643.
\end{thebibliography}
\end{document}
